(1) 设自由费米子Slater波函数为 \(\Psi_F\)，乘 \(A=\prod_{i<j}\mathrm{sgn}(x_i-x_j)\) 得 \(\Psi_B=A\Psi_F\)。两个反对称因子相乘使波函数对称，在粒子重合时仍为零，且 \(|\Psi_B|^2=|\Psi_F|^2\)。奇数 \(N\) 的环可选对称的整数动量填充。

(2) 占据动量 \(k_j=2\pi j/L\)，\(j=-(N-1)/2,\ldots,(N-1)/2\)。因为 \(\sum_jj^2=N(N^2-1)/12\)，得 \(\boxed{E_0=\pi^2\hbar^2N(N^2-1)/(6mL^2)}\)。令 \(n=N/L\) 恒定且 \(N\to\infty\)，单位长度能量 \(\boxed{e=\pi^2\hbar^2n^3/(6m)}\)。

(3) 定义一维压力为 \(P=-\partial_LE_0|_N\)，于是 \(P=\pi^2\hbar^2N(N^2-1)/(3mL^3)\)，连续极限 \(P=\pi^2\hbar^2n^3/(3m)\)。绝热声速 \(c_s^2=(1/m)\partial_nP\)，故 \(\boxed{c_s=\pi\hbar n/m}\)，等于对应自由费米海 \(k_F=\pi n\) 的Fermi速度。

(4) 映射确保对角的粒子密度概率和高阶位置关联与自由费米子相同，但非对角一体密度矩阵含两个位置构型间不同的符号因子，故玻色动量分布不同于Fermi阶跃。有限相互作用需采用Lieb–Liniger模型；偶数粒子的环映射还需处理边界相位。
